%%% FIELD proof-efficient-information
Put \(c=(a,x)\), \(\mu_c=\mu_{ax}(q^0)\), and \(\sigma_c^2=\sigma^2_{ax}(q^0)\).  The support contains at least two distinct values and \(q^0_{c\ell}\geq\eta>0\), by \(\mathrm{def:saturated\text{-}model}\) and \(\mathrm{ass:eta\text{-}interior}\).  Hence
\[
 \sigma_c^2=\sum_{\ell=1}^Lq^0_{c\ell}(y_\ell-\mu_c)^2>0.
 \tag{1}
\]
Define the residual-mean tangent in cell \(c\) by
\[
 m_{c\ell}=\frac{q^0_{c\ell}(y_\ell-\mu_c)}{\sigma_c^2}.
\]
The definitions of \(\mu_c\) and \(\sigma_c^2\) give
\[
 \sum_\ell m_{c\ell}=0,\qquad
 \sum_\ell y_\ell m_{c\ell}=1.                                      \tag{2}
\]
For an arbitrary simplex tangent \(s_c\), set
\[
 b_c=\sum_\ell y_\ell s_{c\ell},\qquad k_c=s_c-b_cm_c.
\]
Because \(\sum_\ell s_{c\ell}=0\), equations (2) imply
\[
 \sum_\ell k_{c\ell}=0,\qquad \sum_\ell y_\ell k_{c\ell}=0.          \tag{3}
\]
This decomposition is unique because applying the mean map to \(s_c=bm_c+k_c\) recovers \(b\).  It remains valid when \(L=2\); then the subspace described by (3) is zero-dimensional.

For the multinomial Fisher inner product at \(q^0_c\), equations (2)--(3) yield
\[
 \sum_\ell\frac{m_{c\ell}k_{c\ell}}{q^0_{c\ell}}
 =\frac1{\sigma_c^2}\sum_\ell(y_\ell-\mu_c)k_{c\ell}=0,\qquad
 \sum_\ell\frac{m_{c\ell}^2}{q^0_{c\ell}}=\frac1{\sigma_c^2}.       \tag{4}
\]
Under \(\mathrm{ass:annotation\text{-}ignorability}\), a labelled observation in cell \(c=(a,x)\) occurs with probability \(p_xe_{ax}r_{ax}\), and an unlabelled observation has zero outcome score.  Multiplying (4) by this probability proves both Fisher orthogonality to every zero-mean shape tangent and the efficient cell-mean information
\[
 (I_r)_{cc}=\frac{p_xe_{ax}r_{ax}}{\sigma^2_{ax}(q^0)}.                \tag{5}
\]
Every entry in (5) is positive: \(p_x>0\) and \(e_{ax}>0\) by their stipulated spaces, \(r_{ax}\geq\rho>0\) because \(r\in\mathcal R_B\), and (1) supplies positive variance.  Thus \(I_r\) is invertible.

The arm-\(a\), stratum-\(x\) column of \(A_{\mathrm{map}}\) is \((-1)^{1-a}p_xd_x\).  Consequently direct block multiplication gives
\[
 A_{\mathrm{map}}I_r^{-1}A_{\mathrm{map}}^\top
 =\sum_{x=1}^K\sum_{a=0}^1
 \frac{p_x\sigma^2_{ax}(q^0)}{e_{ax}r_{ax}}d_xd_x^\top
 =\Sigma(r).                                                         \tag{6}
\]
By \(\mathrm{ass:menu\text{-}rank}\), \(A_{\mathrm{map}}=[-F\;F]\) has row rank \(d\).  For \(z\neq0\), therefore \(A_{\mathrm{map}}^\top z\neq0\), and (6) gives
\[
 z^\top\Sigma(r)z=(A_{\mathrm{map}}^\top z)^\top I_r^{-1}
 (A_{\mathrm{map}}^\top z)>0.
\]
Thus \(\Sigma(r)\) is positive definite.  Substituting the definition
\(T_r=I_r^{-1}A_{\mathrm{map}}^\top\Sigma(r)^{-1}\) into (6) now proves
\[
 A_{\mathrm{map}}T_r=I_d,\qquad
 T_r^\top I_rT_r
 =\Sigma(r)^{-1}A_{\mathrm{map}}I_r^{-1}A_{\mathrm{map}}^\top
   \Sigma(r)^{-1}
 =\Sigma(r)^{-1}.                                                     \tag{7}
\]
Finally, applying (2)--(5) cellwise to the defined probability lift gives
\[
 A_{\mathrm{tan}}\mathcal T_rb=A_{\mathrm{map}}T_rb=b,\qquad
 J_r(\mathcal T_rb,\mathcal T_rc)=b^\top\Sigma(r)^{-1}c.              \tag{8}
\]
Equations (4), (5), and (7) are precisely the claimed orthogonality, efficient information, and least-favourable identities.

%%% FIELD proof-nuisance-experiment
Fix \(C_v<\infty\), and let
\[
 \mathcal S_C=\{s(h,v):h\in\mathcal H,\ \lVert v\rVert\leq C_v\}.
\]
This set is compact because \(\mathcal H\) and the closed nuisance ball are compact in finite dimension and \(s(h,v)=S_0h+Wv\) is linear.  By \(\mathrm{thm:local\text{-}chart\text{-}membership}\), which uses \(\mathrm{ass:eta\text{-}interior}\), \(\mathrm{ass:strict\text{-}interior\text{-}slack}\), and \(\mathrm{ass:tied\text{-}baseline}\), for all sufficiently large \(n\), uniformly on this set,
\[
 q_n(h,v)\in(\Delta_{L-1})^{2K},\quad
 P_{p,q_n(h,v)}\in\mathcal M_\eta(p),\quad
 \sqrt n\,\theta_p(q_n(h,v))=h.                                     \tag{1}
\]
This proves the claimed membership for both signs of every bounded tangent.  Also
\[
 A_C:=\sup_{s\in\mathcal S_C}\max_{a,x,\ell}
 \left|\frac{s_{ax\ell}}{q^0_{ax\ell}}\right|<\infty                 \tag{2}
\]
because \(\mathcal S_C\) is compact and \(q^0_{ax\ell}\geq\eta>0\).

Condition on \(\mathscr P_n\).  Then \(\widehat r_n\) is fixed and, by \(\mathrm{def:admissible\text{-}procedures}\), the main-wave records are independent with annotation probabilities \(\widehat r_n\).  Annotation ignorability makes the factors involving \(R\) identical under \(q_n(h,v)\) and \(q^0\).  Thus the one-record likelihood ratio is \(1\) when \(R=0\), and, when \(R=1\) and \(Y^{\mathrm{obs}}=y_\ell\), it is
\[
 \frac{q^0_{AX\ell}+s_{AX\ell}/\sqrt n}{q^0_{AX\ell}}
 =1+\frac{g_s(O)}{\sqrt n}.                                         \tag{3}
\]
For \(|z|\leq1/2\), Taylor's formula with integral remainder gives
\[
 \log(1+z)=z-\frac{z^2}{2}+r(z),\qquad |r(z)|\leq\frac23|z|^3.       \tag{4}
\]
Once \(A_C/\sqrt n\leq1/2\), summing (4) and using (2) therefore gives, simultaneously for \(s\in\mathcal S_C\),
\[
 \log\Lambda_{n,\mathrm{main}}(s)
 =\Delta_n(s)-\frac1{2n}\sum_{i=1}^ng_s(O_i)^2+R_n(s),\qquad
 \sup_{s\in\mathcal S_C}|R_n(s)|\leq\frac{2A_C^3}{3\sqrt n}.       \tag{5}
\]
Under the conditional baseline law, the simplex-tangent identity \(\sum_\ell s_{ax\ell}=0\) and direct summation over observed cells give
\[
 E_0\{g_s(O_i)\mid\mathscr P_n\}=0,\qquad
 E_0\{g_s(O_i)g_t(O_i)\mid\mathscr P_n\}=J_{\widehat r_n}(s,t).     \tag{6}
\]
Choose a basis \(e_1,\ldots,e_D\) of \(\mathcal E\).  Each \(g_{e_j}g_{e_k}\) is bounded because the observed support is finite and \(q^0\geq\eta\).  Conditional Chebyshev inequalities followed by a union bound over the \(D^2\) pairs imply
\[
 \max_{j,k}\left|\frac1n\sum_{i=1}^ng_{e_j}(O_i)g_{e_k}(O_i)
 -J_{\widehat r_n}(e_j,e_k)\right|=o_{P_0}(1).                       \tag{7}
\]
The coordinate image of \(\mathcal S_C\) is compact.  Hence the quadratic form controlled in (7) is uniformly controlled there, and (5) becomes
\[
 \sup_{h\in\mathcal H,\ \lVert v\rVert\leq C_v}
 |\operatorname{rem}_{n,\widehat r_n}(h,v)|=o_{P_0}(1).              \tag{8}
\]
All constants in (2), (5), and (7) are common to the displayed compact fiber.  Together with the assumed locally uniform convergence \(\widehat r_n\to r\), this also proves the asserted local uniformity.

For any fixed \(t_1,\ldots,t_s\in\mathcal E\), the conditional triangular-array summands
\[
 n^{-1/2}(g_{t_1}(O_i),\ldots,g_{t_s}(O_i))
\]
are centered by (6), uniformly bounded by a constant times \(n^{-1/2}\), and have summed covariance converging in probability to \((J_r(t_j,t_k))_{j,k}\).  Their Lindeberg sums are therefore eventually zero for each positive threshold.  The frozen \(\mathrm{lem:finite\text{-}dimensional\text{-}lindeberg\text{-}feller}\) proves conditional convergence to the centered Gaussian vector with that covariance.  Applying it once to the \(D\)-dimensional basis-score vector and then taking linear images proves finite-coordinate convergence uniformly over \(\mathcal S_C\).

We next identify the contrast projection.  Put \(m(s)=B_{\mathrm{mean}}s\).  The cellwise decomposition in \(\mathrm{lem:efficient\text{-}contrast\text{-}information}\) gives, for \(b\in\mathbb R^d\) and \(s\in\mathcal E\),
\[
 \begin{aligned}
 J_r(\mathcal T_rb,s)
 &=(T_rb)^\top I_rm(s)\\
 &=b^\top\Sigma(r)^{-1}A_{\mathrm{map}}m(s)
 =b^\top\Sigma(r)^{-1}A_{\mathrm{tan}}s.                            \tag{9}
 \end{aligned}
\]
Taking \(s=\mathcal T_rc\) and using \(A_{\mathrm{tan}}\mathcal T_r=I_d\) yields
\[
 J_r(\mathcal T_rb,\mathcal T_rc)=b^\top\Sigma(r)^{-1}c.             \tag{10}
\]
Furthermore \(A_{\mathrm{tan}}\Pi_rs=A_{\mathrm{tan}}s\), so \(s-\Pi_rs\in\ker(A_{\mathrm{tan}})\); equation (9) makes this difference \(J_r\)-orthogonal to \(\operatorname{im}(\mathcal T_r)\).  Thus \(\Pi_r\) is the asserted \(J_r\)-orthogonal projection.

The alternative-law limit follows directly.  For \(s=s(h,v)\) and \(t,u\in\mathcal E\), summation under \(q_n(h,v)\) gives, conditional on the pilot,
\[
 E_sg_t(O_i)=\frac{J_{\widehat r_n}(t,s)}{\sqrt n},\qquad
 \operatorname{Cov}_s\{g_t(O_i),g_u(O_i)\}
 =J_{\widehat r_n}(t,u)+O(n^{-1/2}),                                 \tag{11}
\]
where the remainder is uniform for \(s\in\mathcal S_C\), by (2) and finite support.  The centered summands remain uniformly bounded, so the same Lindeberg--Feller argument applies.  With \(t=\mathcal T_rb\), equations (9)--(11), \(A_{\mathrm{tan}}s(h,v)=h\), and \(\widehat r_n\to r\) show that \(\Delta_n\circ\mathcal T_r\) has limiting mean \(\Sigma(r)^{-1}h\) and covariance \(\Sigma(r)^{-1}\), independently of \(v\).  Hence
\[
 U_{n,r}=\Sigma(r)(\Delta_n\circ\mathcal T_r)
 \ \Rightarrow\ N(h,\Sigma(r))                                    \tag{12}
\]
locally uniformly on the fixed nuisance fiber.

It remains to prove the precise risk transfer and account for the pilot.  For a fully labelled pilot record, the likelihood ratio in tangent \(s\) is again \(1+g_s/\sqrt n\), now with \(R=1\).  If \(J_1\) denotes \(J_r\) with all annotation rates equal to one, independence gives
\[
 E_0\Lambda_{n,\mathrm{pilot}}(s)^2
 =\left(1+\frac{J_1(s,s)}n\right)^{m_n}
 \leq\exp\left\{\frac{m_n}{n}J_1(s,s)\right\}=1+o(1),             \tag{13}
\]
uniformly on \(\mathcal S_C\), by \(\mathrm{ass:pilot\text{-}vanishes}\).  Since its mean is one, (13) implies \(\Lambda_{n,\mathrm{pilot}}(s)\to1\) in \(L^2(P_0)\).  Conditional on any pilot realization,
\[
 E_0\{\Lambda_{n,\mathrm{main}}(s)^2\mid\mathscr P_n\}
 =\left(1+\frac{J_{\widehat r_n}(s,s)}n\right)^n
 \leq e^{J_1(s,s)}.                                                  \tag{14}
\]
Equations (13)--(14) give a uniform \(L^2\) bound for every fixed finite vector of full two-wave likelihood ratios, regardless of discontinuities of the pilot selector.

Let \(\nu=\sum_{\alpha=1}^Nw_\alpha\delta_{h_\alpha}\) be a finite prior.  Put \(t_\alpha=\mathcal T_rh_\alpha\).  By (8) of the efficient-information lemma, \(A_{\mathrm{tan}}t_\alpha=h_\alpha\).  Therefore
\[
 t_\alpha-S_0h_\alpha\in\ker(A_{\mathrm{tan}})=\operatorname{im}(W),
\]
so \(t_\alpha=S_0h_\alpha+Wv_\alpha\) for a unique \(v_\alpha\).  Finiteness places all \(v_\alpha\) in one bounded ball, and (1) places both signs of all these least-favourable submodels in \(\mathcal M_\eta(p)\) eventually.

Pointwise minimization over the finite action set shows that the infimum prior regret over all measurable \(\psi_n\) equals
\[
 E_0\min_{j\in[M]}\sum_{\alpha=1}^N
 w_\alpha\ell_{\mathrm{reg}}(h_\alpha,j)\Lambda_n(t_\alpha).        \tag{15}
\]
The integrand is continuous in the finite likelihood vector and is bounded by
\[
 2\sup_{h\in\mathcal H}\lVert h\rVert_\infty
 \sum_\alpha w_\alpha\Lambda_n(t_\alpha).
\]
The \(L^2\) bound (13)--(14) makes this envelope uniformly integrable.  Joint likelihood convergence from (5), (8), and the conditional Gaussian limit therefore passes the expectation in (15) to the limit.  The limiting likelihood coordinates are
\[
 \exp\left\{h_\alpha^\top\Sigma(r)^{-1}Z
 -\tfrac12h_\alpha^\top\Sigma(r)^{-1}h_\alpha\right\},\qquad
 Z\sim N(0,\Sigma(r)),                                               \tag{16}
\]
which are exactly the likelihood ratios of \(N(h_\alpha,\Sigma(r))\).  Thus (15) converges to the finite-prior Gaussian Bayes regret.  This proves the stated risk-transfer conclusion without asserting total-variation convergence of the experiments, and (13) proves that the pilot contribution is asymptotically negligible.

%%% FIELD proof-gaussian-value
Fix \(r\in\mathcal R_B\).  By \(\mathrm{lem:efficient\text{-}contrast\text{-}information}\), \(I_r\) and \(\Sigma(r)\) are positive definite and
\[
 A_{\mathrm{map}}I_r^{-1}A_{\mathrm{map}}^\top=\Sigma(r),\quad
 A_{\mathrm{map}}T_r=I_d,\quad T_r^\top I_rT_r=\Sigma(r)^{-1}.        \tag{1}
\]
To make the full-fiber assertion explicit, consider the Gaussian cell-mean experiment \(Y_G\sim N(b,I_r^{-1})\), with state \(h=A_{\mathrm{map}}b\) and arbitrary nuisance component of \(b\) in \(\ker(A_{\mathrm{map}})\).  The statistic \(Z=A_{\mathrm{map}}Y_G\) has law \(N(h,\Sigma(r))\) by (1).  Hence rules based only on \(Z\) show that the full-fiber value is at most \(L(r)\).

For the reverse inequality restrict the nuisance fiber to \(b=T_rh\), and define
\[
 E_G=Y_G-T_rA_{\mathrm{map}}Y_G.
\]
Equations (1) imply
\[
 E(E_G)=T_rh-T_rh=0,\qquad
 \operatorname{Cov}(E_G,Z)=I_r^{-1}A_{\mathrm{map}}^\top-T_r\Sigma(r)=0. \tag{2}
\]
Joint Gaussianity and (2) make \(E_G\) independent of \(Z\), with a law not depending on \(h\).  Averaging any rule based on \((Z,E_G)\) over this ancillary law produces a randomized rule based on \(Z\) with exactly the same risk at every \(h\).  Thus even this subexperiment has value \(L(r)\), which proves the reverse inequality and hence the full-fiber characterization.

We next prove existence and minimax duality.  Let \(\mu\) be a nondegenerate Gaussian probability measure on \(\mathbb R^d\).  Every \(N(h,\Sigma(r))\) is equivalent to \(\mu\).  Represent a rule by \((\delta_1,\ldots,\delta_M)\in L^\infty(\mu)^M\) satisfying \(\delta_j\geq0\) and \(\sum_j\delta_j=1\).  This set is convex and weak-star closed in the product of unit balls, hence weak-star compact by \(\mathrm{lem:weak\text{-}star\text{-}rule\text{-}compactness}\).  For fixed \(h\), risk is a weak-star continuous linear functional, because the normal density relative to \(\mu\), multiplied by the bounded loss, belongs to \(L^1(\mu)\).  It is also continuous in \(h\).

The probability measures on compact \(\mathcal H\) form a compact convex set in the weak topology.  Expected risk is affine in both the rule and the prior and continuous in the compact-side variables just described.  The hypotheses of \(\mathrm{lem:sion\text{-}minimax}\) therefore hold, giving
\[
 \inf_\delta\sup_{h\in\mathcal H}R_r(h,\delta)
 =\sup_{\nu\in\mathcal P(\mathcal H)}\inf_\delta
   \int R_r(h,\delta)\,d\nu(h),                                    \tag{3}
\]
and weak-star compactness makes the infimum over \(\delta\) attained.

For completeness, finite priors suffice in (3).  Put \(U=\sup_{h\in\mathcal H}\lVert h\rVert_\infty\) and \(R_0=2U\).  Compactness of \(\mathcal R_B\), \(r_{ax}\geq\rho\), and (1) give
\[
 \lambda_0:=\min_{s\in\mathcal R_B}\lambda_{\min}\{\Sigma(s)\}>0. \tag{4}
\]
For probability densities \(p,q\), partition at \(A=\{p\geq q\}\), apply the log-sum inequality on \(A,A^c\), and write \(t=P(A)-Q(A)\).  The resulting Bernoulli divergence has second derivative at least \(4\), and value and first derivative zero at \(t=0\); hence
\[
 d_{\mathrm{TV}}(P,Q)^2\leq\tfrac12D_{\mathrm{KL}}(P\Vert Q).        \tag{5}
\]
For equal-covariance normals, (4)--(5) yield
\[
 d_{\mathrm{TV}}\{N(h,\Sigma(r)),N(g,\Sigma(r))\}
 \leq\frac{\sqrt d}{2\sqrt{\lambda_0}}\lVert h-g\rVert_\infty.    \tag{6}
\]
The loss is bounded by \(R_0\) and satisfies \(\max_j|\ell_{\mathrm{reg}}(h,j)-\ell_{\mathrm{reg}}(g,j)|\leq2\lVert h-g\rVert_\infty\).  Thus every rule obeys the common modulus
\[
 |R_r(h,\delta)-R_r(g,\delta)|
 \leq\left(2+\frac{R_0\sqrt d}{2\sqrt{\lambda_0}}\right)
 \lVert h-g\rVert_\infty.                                          \tag{7}
\]
Pushing any prior to a nearest point of a finite net of \(\mathcal H\) changes every rule's integrated risk by at most (7) times the mesh.  Taking infima and then refining proves that the supremum in (3) equals the supremum over finite priors.

Finally, \(r\mapsto\Sigma(r)\) is continuous on \(\mathcal R_B\), since every \(r_{ax}\geq\rho>0\).  For positive definite \(S,T\), the Gaussian formula
\[
 D_{\mathrm{KL}}\{N(h,T)\Vert N(h,S)\}
 =\tfrac12\{\operatorname{tr}(S^{-1}T)-d-\log\det(S^{-1}T)\}        \tag{8}
\]
is independent of \(h\) and tends uniformly to zero as \(T\to S\) on the compact covariance range.  Equations (5) and (8), followed by the loss bound, show that \(L(r)\) is continuous.  The set \(\mathcal R_B\) is nonempty because \(\rho\mathbf1\in\mathcal R_B\), and it is closed and bounded, hence compact.  Therefore \(L\) attains its minimum, which by definition is \(V^*\).

%%% FIELD proof-higbee-contrast
The defining member properties of \(\mathfrak C_{\mathrm H}\) are exactly the smoothness, singleton interior optimizer, and uniform strong-concavity conditions stated in the first part of the theorem.  At \(\omega_0\), policy coordinates are \((h_1,h_2,h_3)=(h_1,h_2,0)\).  At \(h=(0,0)\), all three coordinates are equal, and hence
\[
 \operatorname*{argmax}_{j\in[3]}h_j=[3].                            \tag{1}
\]
Every relative neighborhood of zero in \([-1,1]^2\) contains points of the respective forms \((t,0)\), \((0,t)\), and \((-t,-t)\), with \(t>0\).  At these points policies \(1\), \(2\), and \(3\), respectively, are uniquely optimal.

Suppose that the injective representation in the claim existed.  Evaluating it at \(h=0\) and using (1) would give
\[
 \operatorname*{argmax}_{\gamma\in\Gamma_{\mathrm H}}W_{\mathrm H}(0,\gamma)
 =\{\iota_{\mathrm H}(1),\iota_{\mathrm H}(2),\iota_{\mathrm H}(3)\}.
\]
Injectivity makes these three elements distinct, contradicting the singleton-optimizer condition in \(\mathfrak C_{\mathrm H}\).  Thus the witness cannot belong to that comparator class while preserving the three labelled choices.

For the observation-model obstruction, the dimension in the saturated tangent chart is, using the witness values \(K=4\), \(L=3\), and \(d=2\),
\[
 2K(L-1)-d=2\cdot4\cdot2-2=14.                                     \tag{2}
\]
By \(\mathrm{lem:efficient\text{-}contrast\text{-}information}\), the \(J_r\)-projection \(\Pi_r\) of these contrast-preserving directions leaves the efficient contrast signal with covariance \(\Sigma(r)\).  At the witness,
\[
 d_1=(1,0)^\top,\quad d_2=(0,1)^\top,\quad d_3=(1,1)^\top,\quad d_4=0.
\]
Only stratum \(3\) contributes off-diagonally, and its coefficient \(\sum_a p_3\sigma^2_{a3}(q^0)/(e_{a3}r_{a3})\) is strictly positive.  Hence \(\Sigma(r)\) is non-diagonal and changes with the annotation rates.  Equations (1)--(2) establish, respectively, the tie-geometry and selective-annotation nuisance obstructions claimed.

%%% FIELD proof-local-lower
Let \(\mathcal A\in\mathfrak A_{\mathrm{det}}\) have deterministic limiting rate \(r\in\mathcal R_B\), and fix a finite prior \(\nu=\sum_{s=1}^Nw_s\delta_{h_s}\) on \(\mathcal H\).  The finite-prior part of \(\mathrm{thm:nuisance\text{-}uniform\text{-}experiment}\) supplies its least-favourable submodels \(q^0+\mathcal T_rh_s/\sqrt n\), places both signs in \(\mathcal M_\eta(p)\) eventually, and places the finitely many corresponding nuisance coordinates in a common finite ball.  By \(\mathrm{thm:local\text{-}chart\text{-}membership}\), that ball is contained in the full local fiber for all sufficiently large \(n\).  Denote its radius by \(C(\nu)<\infty\).

For every \(C_v\geq C(\nu)\), the supremum defining \(\mathcal R_-(\mathcal A)\) contains this finite submodel.  A supremum dominates a prior average, and the Bayes envelope dominates no particular terminal rule.  The finite-prior transfer in \(\mathrm{thm:nuisance\text{-}uniform\text{-}experiment}\) therefore gives
\[
 \mathcal R_-(\mathcal A)\geq B_r(\nu),                              \tag{1}
\]
where \(B_r(\nu)\) is the Bayes regret in \(Z\sim N(h,\Sigma(r))\).  This step uses no convergence or regularity of the terminal selector: pointwise minimization supplies the Bayes envelope for every measurable selector.

Taking the supremum of (1) over finite priors and invoking the finite-prior equality in \(\mathrm{thm:gaussian\text{-}value\text{-}characterization}\) yields
\[
 \mathcal R_-(\mathcal A)\geq\sup_{\nu\ {\rm finite}}B_r(\nu)=L(r). \tag{2}
\]
The order of limits is the advertised one: \(n\to\infty\) is taken for a fixed ball containing the chosen finite prior, and only afterward does \(C_v\to\infty\).  Finally \(r\in\mathcal R_B\), so the definition \(V^*=\min_{s\in\mathcal R_B}L(s)\) gives
\[
 \mathcal R_-(\mathcal A)\geq L(r)\geq V^*.
\]

%%% FIELD proof-certified-computation
Fix an input in \(\mathfrak I_{\mathrm{rep}}\).  Its exact positivity and rank certificates, the definition of \(\Sigma(r)\), and the floor \(r_{ax}\geq\rho\) give effective spectral bounds
\[
 0<\lambda_0:=\lambda_{\min}\{\Sigma(\mathbf1)\}
 \leq\lambda_{\min}\{\Sigma(r)\},\qquad
 \lambda_{\max}\{\Sigma(r)\}\leq
 \Lambda_0:=\lambda_{\max}\{\Sigma(\rho\mathbf1)\}                \tag{1}
\]
for all \(r\in\mathcal R_B\).  These inequalities follow termwise because every summand of \(\Sigma(r)\) is positive semidefinite and \(1\leq1/r_{ax}\leq1/\rho\).

Put \(H_{ax}=p_x\sigma^2_{ax}(q^0)d_xd_x^\top/e_{ax}\).  Since \(r_{ax},s_{ax}\geq\rho\),
\[
 \lVert\Sigma(r)-\Sigma(s)\rVert_F
 \leq C_r\lVert r-s\rVert_\infty,\qquad
 C_r=\rho^{-2}\sum_{a,x}\lVert H_{ax}\rVert_F.                     \tag{2}
\]
If \(S,T\succeq\lambda_0I\) and \(\lVert T-S\rVert_F\leq\lambda_0/2\), set \(E=S^{-1/2}(T-S)S^{-1/2}\).  Then \(\lVert E\rVert_{\mathrm{op}}\leq1/2\).  Diagonalizing \(E\) and using \(u-\log(1+u)\leq u^2\) for \(|u|\leq1/2\) in the Gaussian divergence formula gives
\[
 D_{\mathrm{KL}}\{N(h,T)\Vert N(h,S)\}\leq\tfrac12\lVert E\rVert_F^2.
\]
The entropy-to-total-variation inequality derived in \(\mathrm{thm:gaussian\text{-}value\text{-}characterization}\) therefore gives
\[
 d_{\mathrm{TV}}\{N(h,T),N(h,S)\}
 \leq\frac{\lVert T-S\rVert_F}{2\lambda_0}.                         \tag{3}
\]
Let \(U=\max_{h\in\mathcal H}\lVert h\rVert_\infty\) and \(R_0=2U\).  Regret is bounded by \(R_0\), and the state-risk modulus proved in that theorem is
\[
 |R(h,\delta)-R(g,\delta)|\leq K_H\lVert h-g\rVert_\infty,\qquad
 K_H=2+\frac{R_0\sqrt d}{2\sqrt{\lambda_0}}.                        \tag{4}
\]

We next construct finite feasible design covers.  The map \(P_B\) is the identity on \(\mathcal R_B\) and maps the full box onto \(\mathcal R_B\).  If \(B=\rho\), it is constant.  If \(B>\rho\), write \(D(u)=\sum_{a,x}p_xe_{ax}u_{ax}\).  On \(D(u)\leq B-\rho\) the map is the identity; on the other piece its derivative is
\[
 \frac{B-\rho}{D(u)}I-
 \frac{B-\rho}{D(u)^2}u(p_xe_{ax})_{a,x}^{\top}.
\]
Because \(\sum_{a,x}p_xe_{ax}=1\), \(D(u)>B-\rho\), and \(\lVert u\rVert_\infty\leq1-\rho\), its sup-norm operator norm is at most \(1+(1-\rho)/(B-\rho)\).  The two formulas agree at the boundary, so the same bound is global.  Thus a rational full-box mesh projects to a finite certified cover of \(\mathcal R_B\).  By (2)--(3), covering radius \(\delta_r\) contributes at most
\[
 E_{\mathrm{des}}=\frac{R_0C_r\delta_r}{2\lambda_0}.                 \tag{5}
\]

Choose \(T>U\).  By (1) and a union bound over coordinates,
\[
 \sup_{h,r}P_{N(h,\Sigma(r))}\{Z\notin[-T,T]^d\}
 \leq2d\exp\left\{-\frac{(T-U)^2}{2\Lambda_0}\right\}=:p_T.       \tag{6}
\]
Tile the cube exactly by rational cubes of side \(\gamma\).  If \(f\) is any relevant Gaussian density and \(\bar f\) its coordinatewise cell average, repeated use of the fundamental theorem of calculus and the \(L^1\)-contraction of averaging gives
\[
 \lVert f-\bar f\rVert_1
 \leq\gamma\sum_{j=1}^d\lVert\partial_jf\rVert_1
 \leq\frac{d\gamma}{\sqrt{\lambda_0}}.                             \tag{7}
\]
The last inequality follows because a Gaussian score coordinate has first absolute moment at most the square root of its variance, and \((\Sigma^{-1})_{jj}\leq1/\lambda_0\).  Averaging any rule on each cube and fixing one action outside changes its risk by at most
\[
 E_{\mathrm{obs}}=R_0\left(p_T+\frac{d\gamma}{\sqrt{\lambda_0}}\right). \tag{8}
\]

At a projected design-grid point \(r\), take a rational \(\xi\)-net of the represented rational polytope \(\mathcal H\), and let \(v(r)\) be the finite zero-sum game value for that state net and the rectangle-constant rules just constructed.  Averaging an arbitrary Gaussian rule proves the first inequality below; applying (4) from a nearest net point to an optimal finite rule proves the second:
\[
 v(r)-E_{\mathrm{obs}}\leq L(r)\leq v(r)+K_H\xi.                    \tag{9}
\]

Every rectangle probability is evaluated by \(\operatorname{GRect}\).  Positive definiteness in (1) makes the algebraic Cholesky conditioning valid.  After tails of total mass at most \(\zeta_{\mathrm{tail}}\) are removed, the iterated integral is over a compact rectangle; all Gaussian derivatives there have finite effective algebraic bounds.  Adaptive subdivision therefore reaches quadrature width \(\zeta_{\mathrm{quad}}\), and increasing interval precision reaches outward-rounding width \(\zeta_{\mathrm{round}}\).  This proves termination and the asserted sum-of-budgets width for each call.

The resulting finite game is a linear program with rational interval coefficients.  Its variables are, explicitly, the action probabilities \(\delta_{jm}\geq0\) for policy \(j\) on observation rectangle \(m\), subject to \(\sum_{j=1}^M\delta_{jm}=1\), together with the epigraph risk variable; the adversary ranges over the finite state net.  Rational primal and dual feasible vectors give outward lower and upper values by weak duality.  Enumeration of rational feasible vectors, mesh refinement, and interval-precision refinement reduce the certified gap below any chosen \(E_{\mathrm{LP}}>0\).  If the maximum rowwise sum of rectangle-probability enclosure widths is \(e_P\), its risk effect is at most \(E_P=R_0e_P\).

Let \(v_{\mathrm{grid}}\) be the smallest certified finite-game value on the design grid.  Combining (5) and (9), in both approximation directions, gives the global enclosure
\[
 \begin{split}
 v_{\mathrm{grid}}-E_{\mathrm{obs}}-E_{\mathrm{des}}-E_{\mathrm{LP}}-E_P
 &\leq V^*\\
 &\leq v_{\mathrm{grid}}+K_H\xi+E_{\mathrm{LP}}+E_P.                \tag{10}
 \end{split}
\]
Every error in (10) is effective and tends to zero as its enumerated mesh or tolerance is refined.  Thus refinement terminates once the width is at most the requested \(\varepsilon\).  The minimizing grid point is feasible by construction, and its certified finite-game rule, with (8)--(10), is \(\varepsilon\)-optimal after allocating the total error budget to be at most \(\varepsilon\).

Finally, cover \(\mathcal R_B\) by the finitely many projected design cells.  For each cell subtract the modulus (5), evaluated at its certified radius, from its finite-game lower bound, and discard it only when that lower bound exceeds the global upper endpoint in (10).  A true minimizer cannot be discarded, so the surviving finite union contains every minimizer.  The construction uses the exact representations required by \(\mathfrak I_{\mathrm{rep}}\); without represented reals or a finite encoding of the compact state set, none of the enumerations is asserted to terminate.

%%% FIELD proof-witness-separation
At \(\omega_0\), the loading vectors are
\[
 d_1=(1,0)^\top,\quad d_2=(0,1)^\top,\quad d_3=(1,1)^\top,\quad d_4=0.
\]
Because \(e_{ax}=1/2\) and every cell variance is \(1/4\), averaging the two arm rates within a stratum preserves annotation cost, while strict convexity of \(t\mapsto1/t\) weakly decreases that stratum's covariance coefficient, with equality only for equal arm rates.  For symmetric useful-stratum rates write
\[
 a=\frac{0.03}{r_1},\qquad b=\frac{0.30}{r_2},\qquad c=\frac{0.25}{r_3}.
\]
Then
\[
 \Sigma(r)=\begin{pmatrix}a+c&c\\c&b+c\end{pmatrix},\qquad
 G_{\omega_0}(r)=\max\{a+c,b+c,a+b\}.                               \tag{1}
\]

Set \(\alpha=9/634\) and
\[
 H_\alpha(r)=\alpha(a+b)+(1-\alpha)(b+c).
\]
This convex combination is at most (1).  As a function of the six useful arm rates, \(H_\alpha\) is strictly convex and strictly decreasing in each rate.  Thus all budget not spent at the forced floor in stratum \(4\) is used.  The Lagrange equations for the remaining linear budget constraint give the unique ratios
\[
 \frac{r_{1G}}{r_{3G}}=\sqrt{\frac\alpha{1-\alpha}}=\frac3{25},\qquad
 \frac{r_{2G}}{r_{3G}}=\frac1{\sqrt{1-\alpha}}=\frac{\sqrt{634}}{25}=:s.
\]
Since the remaining budget is \(11/20-(21/50)(1/20)=529/1000\), the solution is
\[
 r_{3G}=\frac{529/1000}{317/1250+(3/10)s},\quad
 r_{2G}=sr_{3G},\quad r_{1G}=\frac3{25}r_{3G},\quad r_{4G}=\frac1{20}. \tag{2}
\]
Exact comparison after squaring the positive algebraic quantities shows all three useful rates in (2) lie in \((1/20,1)\).  At (2), \(a=c<b\), and hence \(H_\alpha(r_G)=a+b=b+c=G_{\omega_0}(r_G)\).  For every feasible \(r\),
\[
 G_{\omega_0}(r)\geq H_\alpha(r)\geq H_\alpha(r_G)=G_{\omega_0}(r_G). \tag{3}
\]
Strict convexity, its equality condition under arm symmetrization, and strict decrease in useful rates force equality in (3) only at (2), with both agreement-stratum rates at their floor.  Therefore \(\mathcal G(\omega_0)=\{r_G\}\).

For the lower certificate take the prior atoms and weights
\[
 h^0=(-5/8,-5/8),\quad h^1=(0,5/8),\quad h^2=(5/8,0),\qquad
 (w_0,w_1,w_2)=(13/40,7/20,13/40).                                  \tag{4}
\]
At atom \(h^k\), action \(k\) has zero regret and either other action has regret \(5/8\).  With
\[
 \beta_k=\Sigma(r_G)^{-1}h^k,\qquad
 b_k=\log w_k-\tfrac12(h^k)^\top\Sigma(r_G)^{-1}h^k,
\]
the Bayes rule selects the largest affine score \(\beta_k^\top z+b_k\).  If \(P_{kk}\) denotes its correct-selection probability under atom \(k\), its Bayes risk is exactly
\[
 B_G=\frac58\left(1-\sum_{k=0}^2w_kP_{kk}\right)\leq L_{\omega_0}(r_G). \tag{5}
\]
Each \(P_{kk}\) is a bivariate Gaussian probability of two affine inequalities.  Differentiating the bivariate normal distribution in its correlation gives
\[
 \Phi_2(u,v;q)=\Phi(u)\Phi(v)+\int_0^q\phi_2(u,v;t)\,dt.              \tag{6}
\]
For \(|t|\leq q_*<1\), diagonalization of its quadratic form and \(\sup_z z^2e^{-z^2/2}=2/e\) give
\[
 |\partial_t\phi_2(u,v;t)|\leq
 \frac{q_*+2/e}{2\pi(1-q_*^2)^{3/2}}.                               \tag{7}
\]
Thus \(N\)-panel midpoint quadrature for (6) has outward error at most \(q_*^2\) times (7) divided by \(4N\).  Applying the certified arithmetic of \(\mathrm{thm:certified\text{-}computation}\) with \(N=512\) and \(96\)-bit outward rounding to the algebraic design (2), equations (4)--(7) return the strict enclosure
\[
 B_G>0.1965.                                                         \tag{8}
\]

For the feasible alternative, exact rational arithmetic gives
\[
 \sum_xp_xr^0_{\mathrm{alt},x}=\frac{11}{20},\qquad
 \Sigma(r^0_{\mathrm{alt}})=
 \begin{pmatrix}433/1205&125/482\\125/482&8695/14942\end{pmatrix}. \tag{9}
\]
Use the feasible rule selecting the largest of \(z_1,z_2,0\).  With \(v_\Sigma=\Sigma_{11}+\Sigma_{22}-2\Sigma_{12}\), its selection probabilities are
\[
 \begin{aligned}
 P_1(h)&=\Phi_2\!\left(\frac{h_1}{\sqrt{\Sigma_{11}}},
 \frac{h_1-h_2}{\sqrt{v_\Sigma}};
 \frac{\Sigma_{11}-\Sigma_{12}}{\sqrt{\Sigma_{11}v_\Sigma}}\right),\\
 P_2(h)&=\Phi_2\!\left(\frac{h_2}{\sqrt{\Sigma_{22}}},
 \frac{h_2-h_1}{\sqrt{v_\Sigma}};
 \frac{\Sigma_{22}-\Sigma_{12}}{\sqrt{\Sigma_{22}v_\Sigma}}\right),
 \end{aligned}                                                       \tag{10}
\]
and its risk is
\[
 R(h)=\max(0,h_1,h_2)-h_1P_1(h)-h_2P_2(h).                           \tag{11}
\]
The common state-risk modulus on a square of sup-norm half-width \(w\) is
\[
 \left(2+\frac{\sqrt2}{\sqrt{\lambda_{\min}\{\Sigma(r^0_{\mathrm{alt}})\}}}\right)w. \tag{12}
\]
Using (6)--(7) with \(N=64\), a directed dyadic cover of \([-1,1]^2\) with \(74107\) accepted leaves and maximum depth \(10\), and adding (12) on every leaf, the outward certificate gives
\[
 \sup_{h\in[-1,1]^2}R(h)<0.184529576390235<0.18453.                 \tag{13}
\]
There is no observation truncation in (10).  Equations (5), (8), and (13) prove
\[
 L_{\omega_0}(r_G)>0.1965,\qquad
 L_{\omega_0}(r_{\mathrm{alt}}(\omega_0))<0.18453.
\]
The difference is strictly greater than \(0.1965-0.18453=0.01197\).  Since the displayed alternative is feasible, no minimizer of \(L_{\omega_0}\) can equal the unique \(G_{\omega_0}\)-minimizer.

%%% FIELD proof-open-neighborhood
Fix \(\omega\in\mathcal U_{10^{-6}}\).  We verify all numerical inequalities on the closed coordinate box obtained by replacing the strict radii by weak ones; restriction to the admissible relative open subset then proves the claim.  Write
\[
 \sigma_{ax}=\sqrt{\sigma^2_{ax}(q^{0,\omega})},\quad
 S_x=\sigma_{0x}+\sigma_{1x},\quad A_x=p_x^\omega S_x^2,\quad
 t_x=\frac{r_{0x}+r_{1x}}2,
\]
and
\[
 b_x(r)=2p_x^\omega\sum_{a=0}^1\frac{\sigma_{ax}^2}{r_{ax}}.
\]
All denominators are positive because \(r_{ax}\geq\rho>0\), while \(S_x>0\) follows from eta-interiority and the distinct support values.  Cellwise Cauchy--Schwarz gives
\[
 S_x^2
 \leq(r_{0x}+r_{1x})\sum_{a=0}^1\frac{\sigma_{ax}^2}{r_{ax}},\qquad
 b_x(r)\geq\frac{A_x}{t_x},                                        \tag{1}
\]
with equality if and only if
\[
 r_{ax}=\frac{2t_x\sigma_{ax}}{S_x},\qquad a\in\{0,1\}.             \tag{2}
\]
The fixed loading vectors imply
\[
 G_\omega(r)=\max\{b_1(r)+b_3(r),\ b_2(r)+b_3(r),\ b_1(r)+b_2(r)\}. \tag{3}
\]

Define
\[
 \begin{gathered}
 C=B^\omega-\rho p_4^\omega,\qquad
 C_{13}=p_1^\omega A_1+p_3^\omega A_3,\qquad C_2=p_2^\omega A_2,\\
 z=\frac{(\sqrt{C_{13}}+\sqrt{C_2})\sqrt{C_{13}}}{C},\qquad
 y=\frac{(\sqrt{C_{13}}+\sqrt{C_2})\sqrt{C_2}}{C},\\
 t_1^\star=A_1/z,\qquad t_2^\star=A_2/y,\qquad t_3^\star=A_3/z,\qquad
 t_4^\star=\rho.                                                     \tag{4}
 \end{gathered}
\]
All quantities divided by in (4) are positive on the closed box.  Directed interval evaluation at radius \(10^{-6}\), using rational endpoints and outward square-root rounding, verifies
\[
 C_{13}<C_2,\qquad
 \rho<\frac{2t_x^\star\sigma_{ax}}{S_x}<1
 \quad(x\leq3,\ a\in\{0,1\}).                                     \tag{5}
\]
Define \(r_G(\omega)\) by (2) with \(t_x=t_x^\star\) for \(x\leq3\), and put \(r_{a4}=\rho\).  Equations (4)--(5) give the box constraints, and
\[
 \sum_{x,a}p_x^\omega e_{ax}r_{G,ax}(\omega)
 =\rho p_4^\omega+\frac{C_{13}}z+\frac{C_2}y
 =\rho p_4^\omega+C=B^\omega.                                     \tag{6}
\]
Thus this design is feasible.

Set \(\alpha=p_1^\omega A_1/C_{13}\in(0,1)\).  From (1) and (3), every feasible design satisfies
\[
 \begin{aligned}
 G_\omega(r)
 &\geq\alpha\{b_1(r)+b_2(r)\}+(1-\alpha)\{b_2(r)+b_3(r)\}\\
 &\geq\frac{\alpha A_1}{t_1}+\frac{A_2}{t_2}
       +\frac{(1-\alpha)A_3}{t_3}.                                  \tag{7}
 \end{aligned}
\]
Feasibility and \(t_4\geq\rho\) imply
\[
 p_1^\omega t_1+p_2^\omega t_2+p_3^\omega t_3\leq C.               \tag{8}
\]
Cauchy--Schwarz applied to (7), followed by (8), yields
\[
 \begin{aligned}
 \frac{\alpha A_1}{t_1}+\frac{A_2}{t_2}+
 \frac{(1-\alpha)A_3}{t_3}
 &\geq\frac{\left\{\sqrt{\alpha A_1p_1^\omega}+
 \sqrt{A_2p_2^\omega}+\sqrt{(1-\alpha)A_3p_3^\omega}\right\}^2}{C}\\
 &=\frac{(\sqrt{C_{13}}+\sqrt{C_2})^2}{C}=z+y.                     \tag{9}
 \end{aligned}
\]
Here \(\sqrt{\alpha A_1p_1^\omega}=p_1^\omega A_1/\sqrt{C_{13}}\) and the analogous stratum-\(3\) identity prove the second line.  At \(r_G(\omega)\), equality holds in (1), and
\[
 b_1(r_G)=z,\qquad b_2(r_G)=y,\qquad b_3(r_G)=z.
\]
By (5), \(z<y\), so (3) equals \(z+y\).  Equality in (9) uniquely fixes \(t_1,t_2,t_3\) at (4), and equality in (8) fixes \(t_4=\rho\).  Since each fourth-stratum arm rate is at least \(\rho\), their average can equal \(\rho\) only when both equal \(\rho\); equality in (1) fixes both arm rates in each informative stratum via (2).  Therefore
\[
 \mathcal G(\omega)=\{r_G(\omega)\}.                               \tag{10}
\]

For the alternative put
\[
 D_\omega=\sum_xp_x^\omega r^0_{\mathrm{alt},x}-\rho>0,\qquad
 \kappa_\omega=\min\left\{1,\frac{B^\omega-\rho}{D_\omega}\right\}.
\]
The definition of \(P_\omega\) gives
\[
 r_{\mathrm{alt},ax}(\omega)=\rho+\kappa_\omega
 (r^0_{\mathrm{alt},x}-\rho).                                      \tag{11}
\]
If the raw vector is below budget, \(\kappa_\omega=1\) and its cost is at most \(B^\omega\).  If it is above budget, \(\kappa_\omega<1\) and its cost is exactly \(\rho+\kappa_\omega D_\omega=B^\omega\).  Thus (11) is feasible in both branches.  At \(\omega_0\), exact arithmetic gives raw cost \(11/20\), so the projection is the identity there; no unasserted positive budget slack is used.

The outward computation from \(\mathrm{thm:certified\text{-}computation}\), split over the two branches in (11), gives on the entire closed box
\[
 \begin{aligned}
 \lVert\Sigma_\omega(r_G(\omega))-\Sigma_{\omega_0}(r_G(\omega_0))\rVert_F
 &\leq4.0459382\times10^{-5},\\
 \lVert\Sigma_\omega(r_{\mathrm{alt}}(\omega))-
 \Sigma_{\omega_0}(r^0_{\mathrm{alt}})\rVert_F
 &\leq1.2428859\times10^{-5}.                                      \tag{12}
 \end{aligned}
\]
All matrices in these enclosures are positive definite by positive cell weights and full row rank.  Propagating (12) through the center dual prior certificate and center primal rule from \(\mathrm{thm:witness\text{-}separation}\) gives outward total-variation bounds \(7.058912\times10^{-5}\) and \(3.298190\times10^{-5}\), respectively.  The dual-prior loss is bounded by \(5/8\), and regret on \([-1,1]^2\) is bounded by \(2\).  Hence the adverse risk changes are bounded by
\[
 \frac58(7.058912\times10^{-5})=4.4118200\times10^{-5},\qquad
 2(3.298190\times10^{-5})=6.5963800\times10^{-5}.                   \tag{13}
\]
Using the strict center certificates and weak duality or primal feasibility as appropriate,
\[
 \begin{aligned}
 L_\omega(r_G(\omega))
 &>0.1965-0.0000441182=0.1964558818>0.1964558,\\
 L_\omega(r_{\mathrm{alt}}(\omega))
 &<0.18453+0.0000659638=0.1845959638<0.1845960.                     \tag{14}
 \end{aligned}
\]
Since \(0.1964558>0.1845960\), equations (10)--(14) prove the asserted strict separation.

Finally, the neighborhood is nonempty within the stated causal model.  Every allowed standard deviation is realized on \(\{-1,0,1\}\) by
\[
 q^{0,\omega}_{ax}=(\sigma_{ax}^2/2,\ 1-\sigma_{ax}^2,\ \sigma_{ax}^2/2).
\]
These probabilities sum to one, have mean zero and variance \(\sigma_{ax}^2\), and outward endpoint comparison on \(|\sigma_{ax}-1/2|\leq10^{-6}\) puts every coordinate strictly above \(\eta=1/10\).  The zero means imply the required policy tie.  Thus the displayed relative open set contains admissible perturbations.

%%% FIELD proof-binary-reduction
When \(M=2\), the definition \(d=M-1\) gives \(d=1\), and the game in \(\mathrm{def:gaussian\text{-}game}\) observes \(N(h,\Sigma(r))\).  By the definition of the binary benchmark,
\[
 L(r)=L_{\sqrt{\Sigma(r)}}(\mathcal H).                              \tag{1}
\]
If \(Z=h+\sigma\varepsilon\), put \(g=h/\sigma\).  Then \(Z/\sigma\sim N(g,1)\), and positive homogeneity of binary regret gives
\[
 \ell_{\mathrm{reg}}(h,j)=\sigma\ell_{\mathrm{reg}}(g,j).
\]
The maps \(h\leftrightarrow g\) and \(Z\leftrightarrow Z/\sigma\) are bijective because \(\sigma>0\).  Taking the same infimum over rules and supremum over the reparameterized state set therefore proves
\[
 L_\sigma(\mathcal H)=\sigma L_1(\mathcal H/\sigma).                \tag{2}
\]

Now let \(0<\sigma_1\leq\sigma_2\).  Given \(Z_1\sim N(h,\sigma_1^2)\), generate independently \(E\sim N(0,\sigma_2^2-\sigma_1^2)\).  Then \(Z_1+E\sim N(h,\sigma_2^2)\).  Consequently the smaller-noise experiment can reproduce, by this Markov kernel, the risk vector of every rule in the larger-noise experiment.  Taking minimax values yields
\[
 L_{\sigma_1}(\mathcal H)\leq L_{\sigma_2}(\mathcal H).            \tag{3}
\]
Thus a design minimizing the sole contrast variance also minimizes the binary regret value in (1).  The simulation argument proves only the weak inequality (3), so it supplies neither a converse nor strict monotonicity without an additional condition on \(\mathcal H\).
